\documentclass[11pt,a4paper]{article}
\usepackage[margin=2.4cm]{geometry}
\usepackage{booktabs}
\usepackage{amsmath}
\usepackage[T1]{fontenc}
\usepackage{lmodern}
\usepackage{microtype}
\usepackage[hidelinks]{hyperref}
\newcommand{\drop}{\ensuremath{\Delta\text{ipTM}}}
\setlength{\tabcolsep}{5pt}
\renewcommand{\arraystretch}{1.15}

\title{A binder for the KRAS G12D neoepitope on HLA-A*11:01:\\
what twenty designs rest on, and what they do not}
\author{Anthropic $\times$ Adaptyv Protein Design Competition, Challenge 5}
\date{6 November 2026}

\begin{document}
\maketitle

\begin{abstract}
\noindent
The challenge asks for a protein that binds the KRAS G12D peptide \texttt{VVVGADGVGK}
presented on HLA-A*11:01 and not the wild-type \texttt{VVVGAGGVGK} on the same allele, ranked
on mutation specificity first. We report that the first tier is not reachable with the tools
available, and we report it as a measurement rather than as a difficulty. The one atom that
chemically distinguishes the two targets, the G12D carboxylate oxygen OD1, has a
\textbf{0.013\%} escape cone at 8\,\AA\ reach --- 34 of 262{,}144 rays --- so no designed
side chain can touch it. The discrimination that remains is indirect: a 0.85/0.73/0.48\,\AA\
backbone shift at peptide positions 6/7/8, against a 0.337\,\AA\ superposition residual, which
is below the coordinate error of the co-fold predictors used. Measured on the one receptor known
to discriminate these peptides, the predictor's wild-type-versus-mutant sign is not merely
unresolved but \emph{inverted}, by $+0.0027 \pm 0.0006$ over eight independent invocations.

We therefore designed for the geometry that is reachable --- the conformational bulge at peptide
positions 5--8 --- and ranked on peptide-\emph{dependence}: how much a predicted interface
degrades when the peptide is removed from the groove. Calibrated against the 7OW6 receptor over
ten independent invocations the bar is $+0.106$. Of the 78 screened designs, measured twice
independently at five diffusion samples per condition, \textbf{0 are above that bar, 4 are tied
with it, and 74 are below}; 14 of 78 clear zero. Twenty designs are submitted, which is the cap.
We state the funnel with every denominator, declare which filters were fixed after seeing
results, and list what is not established --- including that only one co-fold predictor was
used, so the project's own multi-method requirement is unmet.
\end{abstract}

\section{The target, measured before anything was designed}

The inputs were frozen from three depositions of the same system: \href{https://www.rcsb.org/structure/7OW6}{7OW6}
(G12D ternary complex, HLA heavy chain A, $\beta_2$m B, peptide C, TCR D/E),
\href{https://www.rcsb.org/structure/7OW5}{7OW5} (wild type), and
\href{https://www.rcsb.org/structure/7OW4}{7OW4} (receptor-free). The TCR chains were deleted for
the generation target. One heavy chain and one $\beta_2$m serve all three co-fold conditions,
which differ only in the peptide: the two depositions disagree by a leading construct residue
(276 against 275), and carrying that difference into the comparison would have placed a sequence
change inside the exact contrast selectivity is read from.

\begin{table}[h]
\centering
\caption{Decisive target measurements, all taken before any design was generated.}
\small
\begin{tabular}{p{5.8cm}p{4.3cm}r}
\toprule
Quantity & Condition & Value \\
\midrule
Peptide share of complex accessible surface & --- & 1.2\% (239.5 of 19{,}612\,\AA$^2$) \\
Position-6 side-chain SASA (mutant) & --- & 8.8\,\AA$^2$ (OD1 6.43, OD2 0.00) \\
Escape cone over OD1 & 8\,\AA\ reach, 1.55\,\AA\ probe & \textbf{0.013\%} (34 / 262{,}144 rays) \\
Escape cone over OD1 & 3 / 4 / 5 / 6 / 7\,\AA & 41.0 / 21.9 / 7.1 / 2.9 / 0.018\% \\
Escape cone, position-7 carbonyl & 8\,\AA\ (probe control) & 39.1\% \\
Escape cone, position-4 carbonyl & 8\,\AA\ (probe control) & 0.00\% \\
TCR contacts to position 6 & 5.5\,\AA & 17, \emph{all} to backbone \\
TCR interface split & 4.5\,\AA & 55 peptide / 220 HLA (share 0.200) \\
Peptide backbone shift, mutant vs.\ wild type & C$\alpha$, positions 6/7/8 & 0.85 / 0.73 / 0.48\,\AA \\
HLA superposition residual & the baseline that shift sits against & 0.337\,\AA \\
\bottomrule
\end{tabular}
\end{table}

Two of these are load-bearing. The 161-fold collapse of the escape cone between 6 and 7\,\AA\
locates the carboxylate at the bottom of a pocket about 6\,\AA\ deep; the probe controls behaved
as a working measurement should, with an exposed neighbouring carbonyl at 39.1\% and a buried one
at 0.00\%. And the 0.85\,\AA\ backbone shift sits only 2.5 times above the residual of the
superposition used to measure it, so a predictor whose coordinate error is of that order cannot
be expected to resolve it.

An earlier version of the escape-cone number was quoted as 0.02\% of 4096 rays. That is a single
surviving ray --- the sampling floor, not a measurement --- and it was the most pessimistic
corner of a probe $\times$ reach grid that had not been run. The grid has since been run: the
conclusion did not change, but the evidence for it did, and the tool now reports the surviving
ray count beside the fraction so a result resting on one ray cannot be quoted again without
that being visible.

\section{What the first ranking tier would require, and why it fails}

Mutation specificity is the competition's first tier. The strongest available test of whether any
co-fold predictor can supply it is to run the test on a molecule whose answer is known: the TCR
in 7OW6, which was crystallised bound to the G12D complex and is specific by construction. If the
predictor scores that receptor higher against the mutant peptide than against the wild type, a
design's score in the same column means something. If not, it does not.

Paired within each of the eight independent invocations that predicted the wild-type condition,
the receptor scores \textbf{wild type above G12D by $+0.0027 \pm 0.0006$}, positive in seven of
the eight. The sign is wrong. This is a stronger statement than ``unresolved'': the scoring
function is anti-correlated with the quantity tier 1 asks for, on the one molecule that can
calibrate it. Every subsequent use of the wild-type column in this work is therefore reported
and never gated.

\section{Method}

Generation used BoltzGen 0.1.4 against the intact complex, in four hotspot variants of 64 designs
each --- \texttt{bulge\_core} (peptide positions 5,6,7), \texttt{bulge\_wide} (4--8),
\texttt{mut\_proximal} (5,6) and \texttt{hub} (7,8). The denominator stays \textbf{256} in every
funnel below, including designs rejected for not being designs.

Verification is Boltz 2.2.1, independent of the generator's own folding, in three conditions per
design: \texttt{complex\_g12d} (binder + HLA + $\beta_2$m + mutant peptide), \texttt{apo\_hla}
(the same without the peptide) and \texttt{complex\_wt}. The ranked quantity is
\begin{equation}
\drop \;=\; \text{ipTM}(\texttt{complex\_g12d}) - \text{ipTM}(\texttt{apo\_hla}),
\end{equation}
the degradation of the predicted interface when the peptide leaves the groove. It is a proxy for
peptide-dependent binding, the competition's second tier, and it is explicitly not affinity.

\subsection{The filters, and when each was fixed}

\begin{table}[h]
\centering
\caption{Gates in order of application. ``After'' means the criterion was set after results were
seen; no statistic derived from such a gate is presented as pre-registered.}
\small
\begin{tabular}{p{2.5cm}p{7.6cm}p{3.1cm}}
\toprule
Gate & Criterion & Fixed \\
\midrule
novelty & $<30\%$ of the design is a verbatim copy of a target or TCR chain, by local alignment & after \\
fold novelty & top FoldSeek hit is neither a known protein at high identity nor the target's own immunoglobulin $\beta$-sandwich class & after \\
engagement & peptide share $\geq 0.200$ and every nominated hotspot touched & criterion before; threshold calibrated on the 7OW6 receptor \\
peptide-dependence & \drop\ measured at five samples per condition & before \\
selectivity & \texttt{complex\_wt} & \textbf{reported, not gated} \\
ranking denominator & the cohort's pooled between-run SD, not a design's own two runs & after \\
\bottomrule
\end{tabular}
\end{table}

The engagement threshold is the one that is genuinely calibrated rather than chosen: 0.200 is the
peptide's share of the 7OW6 receptor's own interface, measured with the same code on the same
cutoff. A design that puts less area on the peptide than the natural receptor does cannot be
reading the peptide.

\subsection{The novelty incident, and a claim withdrawn}

A two-design pilot against a \emph{truncated} target returned $\beta_2$-microglobulin --- a chain
of the target itself --- at 90.5\% and 85.7\% identity. Truncation had left $\beta_2$m's own
interface vacant and the generator filled it with the protein that belongs there. Generation moved
to the intact complex, which reduced the rate from 2 of 2 to 17 of 256 but did not eliminate it,
so novelty is a hard gate applied before selection rather than a precaution.

The gate itself was wrong when first shipped. It used a global alignment normalised by the shorter
sequence, which cannot represent a truncated copy at all: a verbatim 40-residue slice of the
276-residue heavy chain scored \textbf{0.000}, and a 60-residue slice of the TCR $\beta$ chain,
the receptor the rules forbid modifying, also scored 0.000. Both would have passed as novel. The
replacement is the fraction of the design that is a verbatim copy of some target chain by local
alignment, normalised by the design's own length; on the same controls it scores 1.000. It was
chosen on synthetic controls alone, independent of the cohort, and the superseded value is kept in
the output so the change is auditable.

One justification offered for gating novelty first was that ipTM rewards plagiarism threefold, so
a score-first ranking would select inadmissible designs preferentially. Measured on the pilot's
four designs that was true; measured over all 256 it is not. The ratio of medians is
\textbf{0.83$\times$} --- the $\beta_2$m copies score slightly \emph{worse} than admissible
designs, 39 of 188 admissible designs beat the best plagiarised one. The threefold effect was real
but local to target preparation. The claim is withdrawn; novelty still gates first, on
correctness rather than on bias.

\section{The funnel}

\begin{table}[h]
\centering
\caption{Every denominator. The campaign's result is this table, not its first row.}
\small
\begin{tabular}{lrp{8.4cm}}
\toprule
Stage & Surviving & Removed by \\
\midrule
generated (4 variants $\times$ 64) & 256 & --- \\
sequence novelty & 239 & 17 verbatim $\beta_2$m copies, 30.3--92.6\% copied \\
fold novelty & 188 & 16 copies of other known proteins (TM 0.93--0.98), 35 immunoglobulin folds \\
engagement & 78 & 110 below the receptor's 0.200 peptide share, or missing a hotspot \\
co-folded twice, 5 samples & 78 & --- (318 of 318 predictions scored, 0 failures) \\
clears zero on \drop & 14 & 64 whose dependence is indistinguishable from none \\
at or above the receptor's bar & \textbf{4} & 10 below the bar but above zero \\
submitted & \textbf{20} & cap is 20; slots 5--20 are the best of the remainder \\
\bottomrule
\end{tabular}
\end{table}

All twenty slots are used. Four designs are tied with the bar and none is above it, so sixteen of
the twenty are submitted as the best available rather than as designs that met a criterion. That
is stated here rather than left to be inferred from a ranked list.

\section{The measurement of record}

The first ranking used one diffusion sample per condition. A 26-design probe at five samples
\emph{reversed the sign} of two of its top four, which is what a difference of two single draws
should be expected to do, so that screen is retrospectively a pilot. Everything below is five
samples per condition, and the whole 159-spec cohort pass was run twice independently.

\subsection{The bar, over ten independent invocations}

\begin{table}[h]
\centering
\caption{The 7OW6 receptor's \drop, measured ten times from byte-identical specifications. Boltz
seeds once per process, so samples within one invocation share that state; separate invocations
are the only way to observe between-run variation.}
\small
\begin{tabular}{lr}
\toprule
Measurement & \drop \\
\midrule
8 small-batch invocations & $+0.101$, spread 0.004 (span $+0.095$ to $+0.108$) \\
26-spec probe & $+0.104$ \\
159-spec cohort, run 1 & $\mathbf{+0.150}$ \\
159-spec cohort, run 2 & $+0.101$ \\
\midrule
pooled over all ten & $\mathbf{+0.106}$, spread 0.016 \\
nine, excluding run 1's control & $+0.101$ \\
\bottomrule
\end{tabular}
\end{table}

Nine of ten measurements of the same molecule agree and one does not. That run is excluded from
the $+0.101$ figure on a \emph{structural} criterion measured independently of its score, not on
the score itself: its peptide-free prediction scored pTM 0.836 against 0.877 and
receptor-to-HLA pair ipTM 0.710 against 0.786, while pLDDT barely moved and the complex condition
reproduced to three decimal places. It is a docking failure in one condition of one run.

Pairing each design against its own run's control would have been the obvious alternative to
excluding it, and it was tested rather than assumed. Across the twelve designs present in both
runs, design apo scores moved $+0.0006 \pm 0.0150$ while the control's fell 0.045. The shift was
not shared, so pairing would have charged every design for one bad prediction.

\subsection{Two variance components, and the one that is not there}

A set of run means has a variance that \emph{already contains} the within-run sampling error, so
the two subtract rather than add:
\begin{equation}
\sigma_{\text{run}} \;=\; \sqrt{\max\!\left(0,\; \sigma^2_{\text{between}} - \mathrm{SEM}^2_{\text{within}}\right)}.
\end{equation}
On the control, $\sigma_{\text{between}} = 0.004$ against $\mathrm{SEM}_{\text{within}} = 0.007$
gives a negative variance, and a 20{,}000-draw permutation test over the pooled samples puts the
observed spread at the \textbf{4th percentile} of the no-run-effect null. There is no resolved
between-run component for the control, and its single-run standard error is if anything
conservative. An earlier version of this analysis combined the two in quadrature; that
double-counted and is withdrawn.

Designs behave differently: pooled between-run SD \textbf{0.018} against a within-run error of
0.015, leaving a run effect of 0.009. Every per-design verdict divides by the pooled figure over
$\sqrt{\text{runs}}$, because two runs cannot estimate one design's own spread --- the estimate
carries roughly 70\% relative error, and a pair that happens to agree reports 0.000, which would
make a $+0.011$ drop read as significant and a reproducible $+0.112$ read as infinitely so. The
control is excluded from the pool it is the bar for, so the bar cannot widen its own margin.

\subsection{The result}

\begin{table}[h]
\centering
\caption{Peptide-dependence, pooled over two independent invocations.}
\small
\begin{tabular}{lrl}
\toprule
 & \drop & per-run values \\
\midrule
\texttt{mut\_proximal\_16} & $\mathbf{+0.112}$ & $+0.112$, $+0.112$ (between-run SD 0.0001) \\
\texttt{mut\_proximal\_26} & $+0.095$ & $+0.109$, $+0.081$ \\
the 7OW6 receptor, ten runs & $+0.106$ & --- \\
median of 78 & $-0.010$ & --- \\
worst of 78 & $-0.093$ & --- \\
\bottomrule
\end{tabular}
\end{table}

\noindent\textbf{0 designs above the bar, 4 tied, 74 below. 14 of 78 clear zero.} The verdict is
unchanged whether the bar is the ten-run $+0.106$ or the $+0.101$ that excludes the degraded-pose
run, which is the most that can be said in its favour.

\subsection{The bias the replicate was run to test}

An unlucky peptide-free pose \emph{inflates} a drop, because the drop is complex minus apo. The
top of this ranking is therefore enriched for apo-prediction failures as well as for genuine
peptide reading, and the apo term carries 4--10$\times$ the complex term's spread, precisely
because removing the peptide leaves the binder's placement underdetermined. Replication is the
test: artifacts regress toward the middle and real signal does not. The two leaders split on it.
\texttt{mut\_proximal\_16} returned $+0.112$ twice; \texttt{mut\_proximal\_26} fell from $+0.109$
to $+0.081$ and is ranked on the pooled mean.

\section{The finalists against the wild-type peptide}

All twenty submitted designs were co-folded against \texttt{VVVGAGGVGK}. Both peptide conditions
are written into a \emph{single} invocation for every finalist and for the control, and the whole
pass was run twice as independent invocations: 420 predictions, 0 failures. The pairing is the
point. The control's bar is formed inside one invocation, so a design's difference has to be
formed the same way, or the run offset stays in the difference and out of the error bar.

The control's own difference is $+0.0018 \pm 0.0065$ over the two invocations ($+0.0030$,
$+0.0007$). It is still positive --- the predictor prefers the peptide the receptor does not
recognise --- but it is tied with zero at this sample size, so the bar is weak and so is any
ranking against it. Three of twenty finalists sit below it, thirteen are tied, four are above;
seven of twenty are negative in raw sign. Across designs the difference averages
$+0.0090 \pm 0.0058$, a spread over designs rather than over runs. The between-run SD is $0.0092$
against a within-run $0.0093$, leaving a run effect of approximately zero.

\subsection*{The unpaired pass did not reproduce}

An earlier pass over the same twenty designs was completed and then withdrawn before it was
reported. Its difference took the wild-type mean from one invocation and the G12D mean from
others. Measured that way, \texttt{bulge\_core} --- the variant whose interface was placed over
peptide positions 5--8, the conformational bulge itself --- appeared 4 of 4 negative at a mean of
$-0.0149$, with Fisher exact $p = 0.0072$ against 3 of 16 elsewhere. That was the one result in
this campaign that looked mechanistic rather than incidental.

Paired, the same four designs average $+0.0115$ at $p = 0.1011$, because \texttt{bulge\_core\_27}
moves across zero to $+0.0767$.

\begin{table}[h]
\centering
\caption{Grouped by generation variant, paired against the unpaired claim. The grouping was
chosen after reading the per-design table and is post-hoc in both columns.}
\small
\begin{tabular}{llrrr}
\toprule
Variant & Hotspots & $n$ & Mean wt $-$ G12D (paired) & Unpaired claim \\
\midrule
\texttt{mut\_proximal} & 5, 6 & 5 & $-0.0012$ & $-0.0073$ \\
\texttt{bulge\_wide} & 4, 5, 6, 7, 8 & 8 & $+0.0106$ & $+0.0052$ \\
\texttt{bulge\_core} & 5, 6, 7 & 4 & $\mathbf{+0.0115}$ & $\mathbf{-0.0149}$ \\
\texttt{hub} & 7, 8 & 3 & $+0.0183$ & $+0.0209$ \\
\bottomrule
\end{tabular}
\end{table}

The measured order is \texttt{mut\_proximal} $<$ \texttt{bulge\_wide} $<$ \texttt{bulge\_core}
$<$ \texttt{hub}, not the pre-registered one. The only ordinal fact that survives is that
\texttt{hub}, which puts the least area on the peptide, is the most positive group (0 of 3
negative).

\textbf{This is not evidence of selectivity, and no design was re-ranked on it.} The bar's sign
is inverted on the only molecule known to discriminate these peptides; the bar is tied with zero
on two invocations; and the grouping is post-hoc. What the replication bought was the retraction
of a result that looked mechanistic and was an artifact of its own denominator --- the same
lesson the peptide-dependence ranking had already learned, applied one level deeper.


\section{What is not established}

\begin{itemize}\itemsep2pt
\item \textbf{Selectivity is sampled, not engineered.} The only chemically distinguishing atom has
a 0.013\% escape cone; the available discrimination is a sub-\AA ngstr\"om backbone shift below the
predictors' coordinate error.
\item \textbf{Only one co-fold predictor was used.} The standing discipline on this project is
that every load-bearing claim gets an independent second method, and that is not met: the Chai-1
container was never located, so the consensus pass did not run. Every number here is Boltz 2.2.1
scoring structures Boltz 2.2.1 predicted. Replication --- ten invocations of the bar, two of the
cohort --- bounds sampling and run-to-run error but cannot detect a systematic error shared by
every run, and the tier-1 inversion is exactly such a case.
\item \textbf{The wild-type pass was never replicated.} One invocation, so the tier-1 column has
no between-run denominator --- the one thing the drop needed before its own verdicts could be
trusted.
\item \textbf{The aspartate rotamer any design must read has never been observed unliganded.}
7OW4, the receptor-free crystal, has no coordinates for peptide positions 5 and 6; 7OW6 holds the
peptide in a conformation a TCR selected. Every design is built against a receptor-shaped
conformation.
\item \textbf{Affinity is not addressed.} The measured quantity is peptide-\emph{dependence}. No
$K_D$ is predicted and nothing ranks by binding strength; a design tied with the control on
dependence may bind far more weakly.
\item \textbf{No wet-lab measurement of anything.} No $K_D$, no expression, no sensorgram.
\item \textbf{No design was shown to exceed the control on any tier.} That the receptor loses
0.106 ipTM when the peptide is removed establishes the scale a design is judged against. It does
not establish that a design matching that number binds the peptide the way the receptor does.
\end{itemize}

\section*{Reproducing the numbers}

\sloppy
Every figure in this report is read from a committed file, not retyped:
\texttt{data/cofold\_replicates.csv} (per-invocation rows),
\texttt{data/cofold\_spread\_pooled.csv} (the ranking table of record),
\texttt{data/cofold\_wildtype.csv} (the tier-1 column),
\texttt{submission/pmhc\_submission\_evidence.csv} (the twenty submitted designs with their
evidence) and \texttt{docs/experiments/pmhc/verdicts.yaml} (one row per experiment, with its
gate and verdict). The commands that regenerate them are in the project README.

\end{document}
